% ---------------------------------------------------------------------
%  11-conclusion.tex
% ---------------------------------------------------------------------
\section{Conclusion}
\label{sec:conclusion}

Software teams learned long ago that parallel change can only be managed,
and audited, if every line of work has a known starting point and an
explicit moment at which it rejoins the main line. Product management never
got the same discipline. Our tools tell us what is being done and in which
workflow state. They do not tell us from which version of the product the
work started, what else was in flight on the same product, or which
firmware release or hardware revision finally closed it. After more than
seven years of managing hardware, firmware and software products, that is
the gap I kept running into, and HPB is my answer to it.

HPB takes the branch-and-merge model from source code and applies it to the
product. It gives you:
\begin{itemize}[leftmargin=1.5em,itemsep=1pt,topsep=3pt]
  \item a \emph{model} of product lines, baselines, typed branches, merges,
        discards, line forks and release gates, governed by seven rules
        (\cref{sec:onepage,sec:model,app:formal});
  \item a \emph{standard notation}, so that everybody draws the same
        situation the same way (\cref{sec:notation,app:card});
  \item an \emph{operating guide} and \emph{metrics} that fall out of the
        graph (\cref{sec:operating,sec:metrics});
  \item \emph{HPB Markup} and an open-source renderer, so a diagram is a
        few lines of text in any web stack or document
        (\cref{sec:tooling,app:markup}); and
  \item a \emph{prototype} in a product management portal that sits on
        top of existing trackers instead of replacing them
        (\cref{sec:implementation}).
\end{itemize}

HPB does not replace Scrum, Kanban, Stage-Gate or formal configuration
management. It sits above them and makes one thing visible that none of
them shows: the product itself, as it evolves through released and planned
versions, with every piece of work anchored to the version it started from
and the version that closed it.

Whether that visibility improves planning and quality in your organization
is something you can test in two release cycles
(\cref{sec:limitations-evaluate}). The notation, the \LaTeX{} package and
the renderer are open, so the method can be adopted, criticized and
improved in a common visual language.

\paragraph{An invitation.}
Try HPB on one product line for two release cycles. Then share what you
found, the diagram that worked or the counter-example that broke the
method, through the issue tracker of \hpbrepo{} or by email to the
address on the first page. Changes to the method and the notation go
through versioned releases (HPB~1.x), so a diagram drawn today stays
readable tomorrow. Start with one line, draw it honestly, and let the
diagram show you what your current tools have been hiding.
